% from ega0-17-fr.tex
\oldpage[0\textsubscript{IV-1}]{45}
Cela résulte aussitôt de (17.2.1) et (17.2.3).
\end{proof}

Le plus petit nombre $n$ (entier ou $+\infty$) pour lequel les conditions
équivalentes de (17.2.8) sont satisfaites est appelé \emph{dimension
cohomologique globale} (ou simplement \emph{dimension cohomologique}) de $A$
et noté $\operatorname{dim.coh}(A)$.

\begin{proposition}[17.2.9]
\label{0.17.2.9-fr}
Soit $A$ un anneau noethérien. Les conditions suivantes sont équivalentes~:
\begin{enumerate}
  \item[\rm a)] $\operatorname{dim.coh}(A)\leq n$.
  \item[\rm b)] Tout $A$-module de type fini est de dimension injective
  $\leq n$.
  \item[\rm c)] Pour tout couple de $A$-modules de type fini, $M$, $N$, on a
  $\operatorname{Ext}^{n+1}_A(M,N)=0$.
\end{enumerate}
\end{proposition}

\begin{proof}
Cela résulte aussitôt de la définition et de (17.2.4).
\end{proof}

\begin{corollary}[17.2.10]
\label{0.17.2.10-fr}
Si $A$ est un anneau noethérien, on a
\begin{equation}
\label{0.17.2.10.1-fr}
\tag{17.2.10.1}
  \operatorname{dim.coh}(A)
  =\sup_{\mathfrak{m}}
  \bigl(\operatorname{dim.coh}(A_{\mathfrak{m}})\bigr)
\end{equation}
où $\mathfrak{m}$ parcourt le spectre de $A$ (ou l'ensemble des idéaux
maximaux de $A$).
\end{corollary}

\begin{proof}
Cela résulte de (17.2.9) et (17.2.5).
\end{proof}

\begin{proposition}[17.2.11]
\label{0.17.2.11-fr}
Soient $A$ un anneau local noethérien, $k$ son corps résiduel. Pour que
$\operatorname{dim.coh}(A)\leq n$, il est nécessaire que
$\operatorname{Tor}^A_i(k,k)=0$, pour $i>n$, et il suffit que
$\operatorname{Tor}^A_{n+1}(k,k)=0$.
\end{proposition}

\begin{proof}
Compte tenu de (17.2.6), il suffit de prouver que les relations
$\operatorname{dim.coh}(A)\leq n$ et
$\operatorname{dim.proj}_A(k)\leq n$ sont équivalentes. Il est clair que la
première entraîne la seconde par définition. Inversement, si
$\operatorname{dim.proj}_A(k)\leq n$, on a
$\operatorname{Tor}^A_i(M,k)=0$ pour $i>n$ et pour tout $A$-module $M$ en
vertu de (17.2.2, (ii))~; donc $\operatorname{dim.proj}(M)\leq n$, ce qui
prouve la proposition en vertu de (17.2.8).
\end{proof}

\begin{corollary}[17.2.12]
\label{0.17.2.12-fr}
Soit $A$ un anneau noethérien. Pour que l'on ait
$\operatorname{dim.coh}(A)\leq n$, il est nécessaire que, pour tout idéal
maximal $\mathfrak{m}$ de $A$, on ait
$\operatorname{Tor}^{A_{\mathfrak{m}}}_i(A/\mathfrak{m},A/\mathfrak{m})=0$
pour $i>n$, et il suffit que ces relations soient vérifiées pour $i=n+1$.
\end{corollary}

\begin{proof}
Cela résulte aussitôt de (17.2.11) et (17.2.10).
\end{proof}

\begin{proposition}[17.2.13]
\label{0.17.2.13-fr}
Soient $A$, $B$ deux anneaux locaux noethériens,
$\varphi:A\longrightarrow B$ un homomorphisme local faisant de $B$ un
$A$-module plat. Alors on a
\begin{equation}
\label{0.17.2.13.1-fr}
\tag{17.2.13.1}
  \operatorname{dim.coh}(A)\leq\operatorname{dim.coh}(B).
\end{equation}
\end{proposition}

\begin{proof}
Supposons que $\operatorname{dim.coh}(B)=n$ soit finie~; il suffit de prouver
que pour tout couple $(M,N)$ de $A$-modules de type fini, on a
$\operatorname{Tor}^A_i(M,N)=0$ pour $i>n$ (17.2.11). Comme $B$ est un
$A$-module fidèlement plat ($0_{\rm I}$, 6.6.2), il revient au même
($0_{\rm I}$, 6.4.1) de montrer que l'on a
\begin{equation}
\label{0.17.2.13.2-fr}
\tag{17.2.13.2}
  \operatorname{Tor}^A_i(M,N)\otimes_A B=0.
\end{equation}
Or, si $L_\bullet=(L_j)$ est une résolution droite de $M$ par des $A$-modules
libres, il résulte de ce que $B$ est un $A$-module plat que
$L_\bullet\otimes_A B=(L_j\otimes_A B)$ est une résolution droite du
$B$-module $M\otimes_A B$ par des $B$-modules libres~; en outre, on a
$(L_j\otimes_A B)\otimes_B(N\otimes_A B)
=(L_j\otimes_A N)\otimes_A B$, d'où on conclut aussitôt que le premier membre
